\subsection{Modelagem de Ameaças e Risco}

\subsubsection{Modelos Físicos de Pressão Humanitária e Territorial}

Todo impacto civil medido pelo SIRDA fundamenta-se num decaimento métrico, onde a ameaça de um incêndio é inversamente proporcional à distância do ativo. A formulação matricial de cálculo adota o formato:
\begin{equation}
\label{eq:pressao_civil}
\text{Índice de Pressão} = \frac{FRP \times (1 + risco\_inpe)}{\max(distancia_{km}, 0.1)}
\end{equation}

Na \textbf{Pressão na Saúde}, o sistema cruza os polígonos de fogo com os ativos do CNES. Pontos com funcionamento "24h" e categoria "urgência/emergência" são isolados como Unidades Críticas, cujo colapso geraria sobrecarga sistêmica no SUS. Atingindo o índice de pressão 100 (limite heurístico de priorização visual), o cenário da instalação torna-se "Crítico".

O \textbf{Índice de Pressão Territorial} aplica a Equação \ref{eq:pressao_civil} à base de Terras Indígenas (FUNAI). A ameaça territorial é subdividida via intersecções \textit{shapely} como: DENTRO (distância euclidiana zero, configurando invasão de chamas) ou ENTORNO (distância maior que zero, mas sujeita aos efeitos radiais da fumaça).

\subsubsection{Modelagem Topológica e Risco em Cascata (NetworkX)}

No cume analítico (Nível 3), a segurança energética é avaliada em duas frentes complementares: o impacto termal direto e a topologia da rede. A vulnerabilidade térmica de uma instalação elétrica expande o Índice de Pressão civil (Equação \ref{eq:pressao_civil}), ponderando-o pela magnitude estrutural do alvo. A formulação matemática da ameaça imediata sobre uma subestação isolada consolida-se por:
\begin{equation}
\text{Ameaça}_{\text{Subestação}} = \text{Índice de Pressão} \times \ln(capacidade_{MW} + 1)
\end{equation}

Para além do impacto isolado, o sistema avança para simular a propagação do risco sistêmico modelando as Subestações e Linhas de Transmissão do ONS (Operador Nacional do Sistema Elétrico) como um grafo direcional via \textit{NetworkX}. O pré-processamento espacial das geometrias EPSG:4326 recorre a um \textit{Spatial Index} otimizado (\textit{cKDTree}) para extrair uma matriz de adjacência georreferenciada baseada na proximidade dos terminais lógicos da malha elétrica, método reconhecido para a avaliação de vulnerabilidade em topologias elétricas.

A partir do grafo gerado, calcula-se a métrica de \textit{Betweenness Centrality} (Centralidade de Intermediação), quantificando estruturalmente a suscetibilidade de uma subestação induzir blecautes caso seja atingida. Pontuações extremas são normalizadas entre 0-100, em que subestações avaliadas com escore $\geq 80$ (ponto de corte empírico projetado para o topo da distribuição de vulnerabilidade) indicam nós críticos com risco de efeito cascata Severo/Extremo. Caso essas artérias centrais colapsem sob a ameaça física do fogo, a destruição causará interrupções macro-regionais ou bloqueios completos de escoamento elétrico (Figura \ref{fig:diagrama_networkx}).

\begin{figure}[H]
    \centering
    \resizebox{0.9\textwidth}{!}{
    \begin{tikzpicture}[
        font=\sffamily,
        process/.style={rectangle, rounded corners=4pt, draw=gray!50, thick, fill=white, font=\small\bfseries, drop shadow={opacity=0.1, shadow xshift=2pt, shadow yshift=-2pt}, align=center, inner sep=6pt, text width=2.8cm, minimum height=1.3cm},
        source/.style={tape, draw=gray!50, thick, fill=white, font=\small\bfseries, drop shadow={opacity=0.1, shadow xshift=2pt, shadow yshift=-2pt}, align=center, inner sep=6pt, text width=2.8cm, minimum height=1.3cm},
        database/.style={cylinder, shape border rotate=90, cylinder end fill=gray!10, draw=gray!50, thick, fill=white, font=\small\bfseries, drop shadow={opacity=0.1, shadow xshift=2pt, shadow yshift=-2pt}, aspect=0.25, align=center, inner sep=4pt, text width=2.8cm, minimum height=1.3cm},
        graph_node/.style={circle, draw=gray!50, thick, fill=white, font=\small\bfseries, drop shadow={opacity=0.1, shadow xshift=2pt, shadow yshift=-2pt}, align=center, text width=2.5cm, inner sep=2pt, minimum size=3.0cm},
        background_box/.style={rectangle, draw=gray!40, thick, dashed, fill=gray!5, rounded corners=6pt, inner sep=15pt},
        relation/.style={-Stealth, thick, draw=gray!60},
        label_text/.style={font=\scriptsize\ttfamily, fill=white, draw=gray!40, rounded corners=2pt, inner sep=3pt}
    ]

    % --- COLUNA 1: BRONZE ---
    \node[source] (sub) {Subestações\\Brutas};
    \node[source, below=1.5cm of sub] (lt) {Linhas de\\Transmissão};
    
    % --- COLUNA 2: TOPOLOGIA ---
    % Centralizado verticalmente com os dois de entrada, ou descendo
    \node[process, right=3cm of sub] (geo) {Conversão GeoPandas\\(EPSG:4326)};
    \node[process, below=1cm of geo] (ckdtree) {cKDTree\\Spatial Index};
    \node[process, below=1cm of ckdtree] (matriz) {Matriz de\\Adjacência};
    \node[graph_node, below=1cm of matriz] (grafo) {Grafo NetworkX\\Direcional};
    
    % --- COLUNA 3: GOLD ---
    \node[process, right=3cm of grafo] (score) {Score de Risco\\em Cascata};
    \node[database, below=1cm of score] (dim_ons) {Dimensão\\Subestações};
    
    % --- LIMITES (SWIMLANES / COLUNAS) ---
    \coordinate (top_level) at ([yshift=0.8cm]sub.north);
    \coordinate (bottom_level) at ([yshift=-0.5cm]dim_ons.south);
    
    \coordinate (c1_left) at ([xshift=-0.5cm]sub.west);
    \coordinate (c1_right) at ([xshift=0.5cm]sub.east);
    
    \coordinate (c2_left) at ([xshift=-0.5cm]geo.west);
    \coordinate (c2_right) at ([xshift=0.5cm]geo.east);
    
    \coordinate (c3_left) at ([xshift=-0.5cm]score.west);
    \coordinate (c3_right) at ([xshift=0.5cm]score.east);

    % Travando o X na coluna e o Y no level universal
    \coordinate (c1_tl) at (c1_left |- top_level);
    \coordinate (c1_br) at (c1_right |- bottom_level);
    
    \coordinate (c2_tl) at (c2_left |- top_level);
    \coordinate (c2_br) at (c2_right |- bottom_level);
    
    \coordinate (c3_tl) at (c3_left |- top_level);
    \coordinate (c3_br) at (c3_right |- bottom_level);

    % BACKGROUNDS
    \begin{scope}[on background layer]
        % Coluna 1
        \node[background_box, fit=(c1_tl) (c1_br) (lt.west) (lt.east)] (bg_col1) {};
        \node[anchor=north west, font=\bfseries\color{black}, xshift=5pt, yshift=-5pt] at (bg_col1.north west) {Extração ONS};
        
        % Coluna 2
        \node[background_box, fit=(c2_tl) (c2_br) (grafo.west) (grafo.east)] (bg_col2) {};
        \node[anchor=north west, font=\bfseries\color{black}, xshift=5pt, yshift=-5pt] at (bg_col2.north west) {Modelagem Topológica};
        
        % Coluna 3
        \node[background_box, fit=(c3_tl) (c3_br) (dim_ons.west) (dim_ons.east)] (bg_col3) {};
        \node[anchor=north west, font=\bfseries\color{black}, xshift=5pt, yshift=-5pt] at (bg_col3.north west) {Métricas de Centralidade};
    \end{scope}

    % --- ARROWS ---
    % Col 1 -> Col 2
    \draw[relation] (sub.east) -- (geo.west);
    \draw[relation] (lt.east) -- (geo.south west); % Diagonal limpa
    
    % Dentro da Col 2
    \draw[relation] (geo.south) -- (ckdtree.north);
    \draw[relation] (ckdtree.south) -- node[label_text, right=2pt] {Distâncias} (matriz.north);
    \draw[relation] (matriz.south) -- (grafo.north);
    
    % Col 2 -> Col 3
    \draw[relation] (grafo.east) -- node[label_text, above=2pt] {Betweenness} (score.west);
    
    % Dentro da Col 3
    \draw[relation] (score.south) -- (dim_ons.north);

    \end{tikzpicture}
    }
    \caption{Cálculo de Centralidade de Risco via Grafos.}
    \fonte{Elaborada pelos autores.}
    \label{fig:diagrama_networkx}
\end{figure}

